% Figure from MATH 452 notes, section 6. Compile with the notes preamble.
\begin{tikzpicture}[scale=1.0]
  % part of the ball B((x0,y0), r) inside the open set R
  \begin{scope}
    \clip (-1.7,-0.75) rectangle (3.9,2.35);
    \draw[black!55, thin, dashed] (0,0) circle (3.3);
  \end{scope}
  \node[black!70, font=\scriptsize, anchor=north west, inner sep=1pt] at (-1.5,2.2) {$B((x_0,y_0),r) \subseteq R$};
  % the direct increment (h,k)
  \draw[black!45, thin, dashed] (0,0) -- (1.9,1.2);
  % a smaller increment: the L-path shrinks and drags the MVT points to (x0,y0)
  \draw[figB!45, thin] (0,0) -- (0.85,0) -- (0.85,0.54);
  \filldraw[figR!45] (0.3,0) circle (1pt);
  \filldraw[figR!45] (0.85,0.28) circle (1pt);
  % the L-shaped path: first in x, then in y
  \draw[figB, thick] (0,0) -- (1.9,0) node[pos=0.72, above, font=\scriptsize, inner sep=2pt] {$h$};
  \draw[figB, thick] (1.9,0) -- (1.9,1.2) node[pos=0.8, right, font=\scriptsize, inner sep=2.5pt] {$k$};
  % the MVT points
  \filldraw[figR] (0.75,0) circle (1.4pt) node[below, font=\scriptsize, inner sep=2.5pt] {$f_x$};
  \filldraw[figR] (1.9,0.55) circle (1.4pt) node[right, font=\scriptsize, inner sep=2.5pt] {$f_y$};
  % corners
  \filldraw (0,0) circle (1.3pt) node[left, font=\scriptsize, inner sep=2.5pt] {$(x_0,y_0)$};
  \filldraw (1.9,0) circle (1.3pt) node[below, font=\scriptsize, inner sep=2.5pt, xshift=6pt] {$(x_0+h,\,y_0)$};
  \filldraw (1.9,1.2) circle (1.3pt) node[above, font=\scriptsize, inner sep=2.5pt] {$(x_0+h,\,y_0+k)$};
\end{tikzpicture}
